\chapter{健康的选择}\label{ch:g5:healthy-choices}

到目前为止，本书已经带你参观了身体的大部分重要“部门”：
\omterm{def:g4:journey-of-food:tract}{消化道}、肺、\omterm{def:g4:heart-and-blood:heart}{心脏}及其
奔流的\omterm{prop:g4:journey-of-food:absorb}{血液}、\omterm{def:g3:bones-and-muscles:muscle}{肌肉}与\omterm{def:g3:bones-and-muscles:skeleton}{骨骼}，以及发号施令的\omterm{def:g5:brain-in-command:system}{大脑}。
本学年的最后一章提出一个实际的问题：
既然身体有这样复杂的机器，怎样才能把它运转好呢 ——
又由谁来决定？因为大约从你这个年纪开始，越来越多的
决定要由你自己来做。

\section{身体对我们的要求}

\begin{proposition}[身体的日常“契约”]\label{prop:g5:healthy-choices:contract}
把前面各章的要求放在一起看，就会发现一份
日常“契约”：
\begin{itemize}
  \item \emph{\omterm{prop:g1:healthy-habits:needs}{多样的食物}和水}——燃料与建造材料
        （见\cref{prop:g3:balanced-diet:jobs}），
        由\omterm{prop:g4:journey-of-food:absorb}{血液}输送到每个\omterm{def:g5:first-look-at-cells:cell}{细胞}；
  \item \emph{运动}——\omterm{def:g3:bones-and-muscles:muscle}{肌肉}、\omterm{def:g3:bones-and-muscles:skeleton}{骨骼}、\omterm{def:g4:heart-and-blood:heart}{心脏}和肺
        都会因使用而变得更强健
        （见\cref{met:g3:bones-and-muscles:care}，
        \cref{rem:g4:heart-and-blood:rest}）；
  \item \emph{\omterm{prop:g1:healthy-habits:needs}{睡眠}}——\omterm{def:g5:brain-in-command:system}{大脑}的整理归档与身体的修复
        （见\cref{prop:g5:brain-in-command:protect}）；这个年龄的孩子
        约需要 10 小时；
  \item \emph{干净的空气、干净的双手、好好照顾的牙齿}——
        在\cref{ch:g1:healthy-habits}和
        \cref{met:g4:breathing:care}中的防护措施仍在发挥作用。
\end{itemize}
没有新东西——这正是关键。健康就是老规则，
坚持下去。
\end{proposition}
\admitted

\begin{example}[从身体内部看的一天]\label{ex:g5:healthy-choices:day}
一个好的一天，从身体管理的角度看：早餐为
补充夜间消耗的\omterm{prop:g3:balanced-diet:jobs}{能量}（见\cref{ex:g3:balanced-diet:breakfast}）；
上学路上的步行让心肺热身；中午的游戏
其实是伪装成乐趣的锻炼；晚餐的蔬菜为
“新鲜家族”补充成员；而早睡则把
整理归档的时间留给\omterm{def:g5:brain-in-command:system}{大脑}。没有哪一项是
英雄之举；合起来，才构成健康。
\end{example}

\begin{omfigure}
\includegraphics[width=0.6\linewidth]{images/book1/ai/g5-choices-bicycle.jpg}

{\small 同一幅图里体现了“契约”的第二条和第四条：运动
为\omterm{def:g4:heart-and-blood:heart}{心脏}、肺和\omterm{def:g3:bones-and-muscles:muscle}{肌肉}而做的运动——还有一顶
保护“指挥官”\omterm{def:g5:brain-in-command:system}{大脑}的头盔。}
\end{omfigure}

\section{新的选择，新的压力}

\begin{proposition}[屏幕与睡眠]\label{prop:g5:healthy-choices:screens}
屏幕——手机、平板、游戏——既不是毒药
也不是食物；它们是\emph{时间}。它们占去的时间，
都要从别处挤出来，而通常受害的正是
“契约”中的条目：户外运动和\omterm{prop:g1:healthy-habits:needs}{睡眠}。更糟的是，
晚上很晚时的明亮屏幕会告诉\omterm{def:g5:brain-in-command:system}{大脑}“仍然
白天”，使\omterm{prop:g1:healthy-habits:needs}{睡眠}推迟，超出了被占用的时间。
所以问题不在于用不用屏幕，而在于
它们可以占用哪些时段——而晚间的最后一小时
是最不该交给屏幕的时段。
\end{proposition}
\admitted

\begin{example}[十一点实验]\label{ex:g5:healthy-choices:latenight}
不妨把它当作一次公平试验（见\cref{met:g4:what-plants-need:fairtest}）
在自己身上做：睡前一小时关掉屏幕连续一周，
与普通的一周相比。多数人入睡更快，
醒来更轻松，第一周就能感觉早晨更舒服——
对照周本身就是论据。
\end{example}

\begin{proposition}[欺骗身体的物质]\label{prop:g5:healthy-choices:substances}
有些物质恰恰会损害本书中已经
参观过的各个系统：
\begin{itemize}
  \item \emph{烟草烟雾}会弄脏肺部精细的
        肺泡（见\cref{met:g4:breathing:care}），并使
        \omterm{def:g4:heart-and-blood:heart}{心脏}负担加重；它的“推销”靠的是让第一支香烟
        显得“成熟”——而它留住买家的办法是
        \emph{成瘾}\index{成瘾}：身体一旦
        习惯了它，缺少它就会“抗议”；
  \item \emph{酒精}几分钟内就会到达\omterm{def:g5:brain-in-command:system}{大脑}
——\omterm{def:g5:brain-in-command:system}{大脑}这位“指挥官”本身
        （见\cref{ch:g5:brain-in-command}），减慢\omterm{def:g5:brain-in-command:system}{神经}回路，
        使判断变得模糊；年轻的\omterm{def:g5:brain-in-command:system}{大脑}仍在建立
        \omterm{def:g5:brain-in-command:system}{神经}通路，比成年人的\omterm{def:g5:brain-in-command:system}{大脑}更容易受损；
  \item 二者早已被列为对\omterm{def:g5:human-reproduction:pregnancy}{胎儿}有危害
        （见\cref{prop:g5:human-reproduction:care}）——同样的物质，
        同样的理由。
\end{itemize}
\end{proposition}
\admitted

\begin{remark}[为什么“永不开始”胜过“以后再停”]\label{rem:g5:healthy-choices:never}
\omterm{prop:g5:healthy-choices:substances}{成瘾}是一个入口容易、出口艰难的陷阱：
开始只需一个选择，停止则是一场反复
持续数年的斗争。这种不对称是本章最冷峻的事实，
同时也最有用：赢得这场斗争成本最低的时刻
就是它开始之前。递给别人第一支香烟的人，
并不只是递出一支烟——他们递出的是之后的
所有香烟。
\end{remark}

\section{自己做决定}

\begin{method}[自己执行身体的“契约”]\label{met:g5:healthy-choices:run}
从现在起，并且逐年增加，你身体“契约”的守护者
就是你自己：
\begin{enumerate}
  \item 了解“契约”的各项条目——
        \cref{prop:g5:healthy-choices:contract}即为清单；
  \item 时常诚实地检查自己，做法如
        \cref{met:g3:balanced-diet:check}所示——缺了什么，
        哪样超出了它应有的份额？
  \item 每次只改变一件事，并为它留出一个
        公平的一周试验期；
  \item 对于物质，要在需要它的时刻\emph{之前}
        就定好自己的“不”
        （见\cref{rem:g5:healthy-choices:never}）——冷静提前做出的决定，
        比在善意压力下临场拼凑的决定
        更容易坚守；
  \item 还要求助大人：医生、护士、父母和
        老师都是这份契约中站在你这一边的人。
\end{enumerate}
\end{method}

\begin{example}[事先准备好的回答]\label{ex:g5:healthy-choices:answers}
第 4 步所说的提前决定，落到实处就是一句
随时备在口袋里的话：“不了，谢谢——这不是我的菜。”说出口时，
语气轻松、只说一次，通常就能让此事了结；值得交往的朋友
第一次就会接受。这句话准备起来毫无成本，
而偏偏在需要它的那一刻最难组织，
所以现在就把它组织好。
\end{example}

\section{练习}

\begin{exercise}[$\star$]\label{exo:g5:healthy-choices:1}
列出身体“日常契约”中的四项内容，并说明每一项
主要由哪个器官或哪项服务承担。
\end{exercise}

\begin{exercise}[$\star$]\label{exo:g5:healthy-choices:2}
与你同龄的儿童，大约需要多少\omterm{prop:g1:healthy-habits:needs}{睡眠}，
\omterm{prop:g1:healthy-habits:needs}{睡眠}时身体会完成哪两项工作？
\end{exercise}

\begin{exercise}[$\star$]\label{exo:g5:healthy-choices:3}
屏幕时间在什么意义上是“从别处挪来的”？
说出两项最常被挤占的事情。
\end{exercise}

\begin{exercise}[$\star$]\label{exo:g5:healthy-choices:4}
为什么傍晚的最后一个小时
是最不应该拿来对着明亮屏幕的时段？
\end{exercise}

\begin{exercise}[$\star$]\label{exo:g5:healthy-choices:5}
烟草烟雾首先会伤害哪些器官，依据
\cref{prop:g5:healthy-choices:substances}？酒精呢？
\end{exercise}

\begin{exercise}[$\star$]\label{exo:g5:healthy-choices:6}
什么是\omterm{prop:g5:healthy-choices:substances}{成瘾}？为什么“永不开始”
比“以后再停”的代价低得多？
\end{exercise}

\begin{exercise}[$\star$]\label{exo:g5:healthy-choices:7}
提前决定“不”有何意义？结合
\cref{met:g5:healthy-choices:run}中的第 4 步回答。
\end{exercise}

\begin{exercise}[$\star$]\label{exo:g5:healthy-choices:8}
动手试一试：做
\cref{ex:g5:healthy-choices:latenight}中的“十一点”实验，连续一周，
并——诚实地——汇报发生了什么变化。
\end{exercise}

\begin{exercise}[$\star\star$]\label{exo:g5:healthy-choices:9}
设计屏幕实验，使其成为规范的公平试验：
两周之间哪些保持不变，哪一件事
发生改变，你又测量什么？
\end{exercise}

\begin{exercise}[$\star\star$]\label{exo:g5:healthy-choices:10}
“只抽一根，不会有什么伤害。”请结合
\cref{rem:g5:healthy-choices:never}来回答：第一根烟真正
提供的是什么？
\end{exercise}

\begin{exercise}[$\star\star\star$]\label{exo:g5:healthy-choices:11}
请将前面的三章与本章联系起来：
\cref{prop:g4:breathing:exchange}、
\cref{ex:g4:heart-and-blood:effort}和
\cref{prop:g5:brain-in-command:protect}
各自如何解释“契约”中的某一条，或本章的某项警示？
\end{exercise}
